\documentclass[a4paper,14pt]{extarticle}
\usepackage{styles}


\begin{document}

\litSubsubsection{3.2.1 Performance Analysis}

The XOR perceptron problem was selected as the second benchmark to complement the TSP evaluation. XOR is a standard test case for neural network training and genetic algorithm evaluation~\cite{googlebooks_xor_perceptron}. It represents a linearly inseparable classification task that is simple enough to produce consistent reproducible results to demonstrate evolutionary algorithms. An important property of this benchmark is that a single-layer perceptron with two weights can correctly classify at most 2 of 4 XOR input patterns. Both GPU and CPU implementations consistently reach this theoretical ceiling, so execution time (not solution quality) serves as the main metric.

\textbf{Algorithm Configuration.} The primary benchmark uses a population of 100 individuals over 500 generations. The scaling analysis varies population size from 50 to 1000 over 200 generations. A mutation rate of 0.1 is applied to a gene length of 2, encoding perceptron weights $w_0$ and $w_1$. Tournament selection and single-point crossover are used as genetic operators. Each configuration is run 30 times for statistical significance.

\textbf{Implementation Versions.} The CUDA implementation evolved through iterations v0--v5. The final version uses shared memory for read-only training data, pre-allocated device arrays, pinned host buffers, and GPU-side RNG via Xoroshiro128p. Since the XOR fitness function is lightweight compared to TSP, this benchmark isolates transfer and launch overhead more clearly.

\textbf{Performance Metrics.} Four metrics are recorded for each run: wall time (total execution including transfers), compute time (pure kernel execution), transfer time (host--device data movement), and fitness throughput (evaluations per second). Statistical significance is assessed via paired t-test at $p < 0.05$.

\litSubsubsection{3.2.2 Results}

\textbf{Primary Benchmark Results.} Table~\ref{tab:primary_results} shows the primary comparison for population size 100 over 500 generations, averaged across 30 independent runs.

\begin{table}[H]
\centering
\caption{Statistical comparison of GPU vs. CPU performance (30 runs each)}
\label{tab:primary_results}
\begin{tabular}{@{}lcccc@{}}
\toprule
\textbf{Metric} & \textbf{GPU Mean} & \textbf{CPU Mean} & \textbf{p-value} & \textbf{Significant} \\
\midrule
Wall Time (s) & 0.368 $\pm$ 0.044 & 0.556 $\pm$ 0.060 & $< 0.0001$ & Yes \\
Compute Time (s) & 0.248 $\pm$ 0.034 & 0.319 $\pm$ 0.034 & $< 0.0001$ & Yes \\
Transfer Time (s) & 0.107 $\pm$ 0.017 & N/A & N/A & N/A \\
Fitness Evaluations/sec & 201,153 & 156,278 & N/A & N/A \\
\bottomrule
\end{tabular}
\end{table}

The GPU achieves a 1.51$\times$ wall-time speedup over the CPU baseline with high significance ($p < 0.0001$). Host--device data transfer accounts for 29.2\% of total GPU execution time. The GPU evaluates fitness 28.7\% faster than the CPU (201,153 vs.\ 156,278 evaluations per second). Both implementations converge to identical fitness (2 of 4 patterns), which confirms that GPU acceleration preserves solution correctness. The zero-variance fitness result reflects the linear separability ceiling of XOR, not a lack of algorithmic diversity.

\litSubsubsection{3.2.3 Scaling Analysis}

Table~\ref{tab:scaling_results} shows how performance scales across population sizes. The scaling effect is stronger than in the TSP benchmark due to the lightweight per-individual fitness computation. All differences are statistically significant ($p < 0.05$ or better at every population size).

\begin{table}[H]
\centering
\caption{Scaling performance across population sizes (10 runs each, 200 generations)}
\label{tab:scaling_results}
\begin{tabular}{@{}lcccr@{}}
\toprule
\textbf{Pop.} & \textbf{GPU Time (ms)} & \textbf{CPU Time (ms)} & \textbf{Speedup} & \textbf{Overhead \%} \\
\midrule
50 & 138.5 $\pm$ 11.2 & 100.6 $\pm$ 6.1 & 0.73x & 30.5\% \\
100 & 137.1 $\pm$ 21.0 & 195.7 $\pm$ 12.3 & 1.43x & 29.9\% \\
500 & 156.3 $\pm$ 13.6 & 1038.5 $\pm$ 34.5 & 6.64x & 28.8\% \\
1000 & 155.1 $\pm$ 10.4 & 2257.6 $\pm$ 35.3 & 14.56x & 27.6\% \\
\bottomrule
\end{tabular}
\end{table}

The GPU breakeven point occurs at population size $\geq 100$. Below this threshold, transfer overhead dominates and the CPU is faster (0.73$\times$ at population 50). GPU compute time remains nearly constant as population grows, rising only from 138.5\,ms to 155.1\,ms across a 20$\times$ increase in population size. CPU time, by contrast, scales linearly from 100.6\,ms to 2257.6\,ms. Transfer overhead decreases from 30.5\% to 27.6\% as population grows, reflecting better amortization of fixed transfer costs. A 20$\times$ increase in population yields a 19.9$\times$ speedup improvement, which indicates near-linear scaling efficiency.

\begin{figure}[H]
\centering
\begin{subfigure}{0.49\textwidth}
\includegraphics[width=\textwidth]{diagrams/sec3/time_comparison.png}
\caption{Wall time comparison across population sizes}
\label{fig:time_comparison}
\end{subfigure}
\hfill
\begin{subfigure}{0.49\textwidth}
\includegraphics[width=\textwidth]{diagrams/sec3/speedup_vs_population.png}
\caption{Speedup factors vs. population size}
\label{fig:speedup_vs_pop}
\end{subfigure}
\end{figure}

\begin{figure}[H]
\centering
\begin{subfigure}{0.49\textwidth}
\includegraphics[width=\textwidth]{diagrams/sec3/fitness_throughput.png}
\caption{Fitness evaluation throughput}
\label{fig:fitness_throughput}
\end{subfigure}
\hfill
\begin{subfigure}{0.49\textwidth}
\includegraphics[width=\textwidth]{diagrams/sec3/gpu_time_breakdown.png}
\caption{GPU time breakdown (compute vs. transfer)}
\label{fig:gpu_breakdown}
\end{subfigure}
\end{figure}

\begin{figure}[H]
\centering
\includegraphics[width=0.8\textwidth]{diagrams/sec3/scaling_efficiency.png}
\caption{Scaling efficiency analysis}
\label{fig:scaling_efficiency}
\end{figure}

\litSubsubsection{3.2.4 Statistical Significance}

Paired t-tests show significant execution time differences between GPU and CPU across all population sizes from 50 to 2000 ($p < 0.001$, Bonferroni-corrected). All configurations converge to the theoretical fitness ceiling (2 of 4 patterns), with no difference in solution quality between implementations. Cohen's $d$ effect sizes exceed 0.8 (large effects) for populations $\geq 100$, confirming GPU acceleration at realistic scales.

\begin{center}
\begin{tabular}{@{}lcc@{}}
\toprule
\textbf{Effect Size (Cohen's $d$)} & \textbf{Required $N$} & \textbf{Observed Power ($N=30$)} \\
\midrule
Small (0.2) & 393 & 15\% \\
Medium (0.5) & 63 & 52\% \\
Large (0.8) & 25 & 89\% \\
\bottomrule
\end{tabular}
\end{center}

The $N = 30$ design provides 89\% power for large effects. Since the observed speedups at population $\geq 100$ produce large effect sizes ($d > 0.8$), the sample size is sufficient.

\end{document}
    